\section{Fixed sources, primitive families, and fixed moments}
\label{p1:fixed:section}

The boundary estimate will use large families, but its finitely
many small sizes require stronger arithmetic input. We establish
that input first. All numbers of sources and all moment orders
in this section are fixed before \(N\) tends to infinity.

For a family \(I\) of windows, let \(R(I)\) be the kernel
of the sum map on the effective spaces and put \(\rho_{\mathrm{rel}}(I)=\dim R(I)\).
For \(s\) fixed sources \(\boldsymbol x=(x_1,\ldots,x_s)\) and a
new centre \(z\), write
\begin{equation}
 \delta_z=\dim R(\boldsymbol x,z)-\dim R(\boldsymbol x)
          =\dim\operatorname{pr}_z R(\boldsymbol x,z).
 \label{p1:fixed:delta}
\end{equation}
If the prescribed affine events at the sources and target are
jointly compatible, their conditional probability is
\(p_z2^{\delta_z}\); it is zero otherwise. The same formula applies
to homogeneous prescriptions. The proofs on regular windows and their
effective-space extension in Section~\ref{p13:exact:section} establish
the following inputs without using a distributional approximation.

\subsection{Regularizing the coordinate labels}

The following linear-algebra lemma gives control of every short
circuit of coordinate labels, including repeated labels.

\begin{lemma}[Disjoint recovery]
\label{p1:fixed:recovery}
Let \(R\subseteq\F_2^V\), \(n=|V|\), and let
\(\ell_v\in R^*\) be its coordinate forms. For fixed integers
\(d\ge1\) and \(t\ge2\), there is a subspace \(R'\subseteq R\)
with
\begin{equation}
 \codim_R R'\le 2n t^{-1/d}
 \label{p1:fixed:recovery-cost}
\end{equation}
such that no nonzero \(a\in(R')^*\) has \(t\) representations
\[
 a=\sum_{v\in A_i}\ell_v|_{R'},\qquad
 1\le |A_i|\le d,
\]
by pairwise disjoint subsets \(A_1,\ldots,A_t\) of \(V\).
Each nonzero coordinate label occurs at most \(K=t-1\) times.
For \(1\le j\le d\), a nonzero label has at most
\(j!t^{j-1}\) representations as a sum of \(j\) distinct
linearly independent labels.
\end{lemma}

\begin{proof}
Start with \(W=0\subset R^*\). Whenever \(R^*/W\) has a
nonzero class with \(t\) such disjoint recovery sets, record
those sets and adjoin one representative \(\alpha_i\) of the
class to \(W\). Each step increases \(\dim W\) by one. Fix the
entire resulting sequence of \(m\) steps.

Select each coordinate independently with probability
\(u\in(0,1)\), obtaining \(A\subseteq V\). At step \(i\),
declare success if one of its recorded recovery sets is
contained in \(A\). On success,
\(\alpha_i\in\langle\ell_v:v\in A\rangle+W_{i-1}\).
The failure probability is at most \((1-u^d)^t\), since the
recorded sets at that step are disjoint. For each realization,
\[
 m\le |A|+\#\{\hbox{failed steps}\}.
\]
Taking expectations and choosing \(u=t^{-1/d}\) gives
\[
 m\le \frac{nu}{1-(1-u^d)^t}\le2nt^{-1/d}.
\]
The annihilator \(R'=W^\perp\) has the stated codimension
and satisfies the stopping condition.

The multiplicity bound follows by using singleton recovery sets.
To count independent \(j\)-term representations of a nonzero
\(a\), take a maximal disjoint family of their label supports.
It has fewer than \(t\) members. Its union has at most
\(j(t-1)\) labels and meets every representation. Removing
one of these labels reduces to a \((j-1)\)-term independent
representation of a nonzero vector: the removed label cannot
equal \(a\), because the remaining independent labels would
then sum to zero. Induction from \(j=1\) proves the last bound.
\end{proof}

\subsection{Uniformly few targets of large rank loss}

\begin{proposition}[Large rank loss after fixed sources]
\label{p1:fixed:targets}
For every fixed \(s\ge1\) and \(\eta,\varepsilon>0\),
\begin{equation}
 \sup_{x_1,\ldots,x_s\ {\rm distinct}}
 \#\{z\in I_N\setminus\{x_1,\ldots,x_s\}:
                   \delta_z\ge\eta L\}
       \ll_{s,C,\eta,\varepsilon}N^\varepsilon .
 \label{p1:fixed:targets-bound}
\end{equation}
The sources may be dependent or overlapping.
\end{proposition}

\noindent The complete proof is given in Section~\ref{p14:comp:fixed-targets}.


\begin{proposition}[Conditional sums and propagation]
\label{p1:fixed:propagation}
For every fixed \(s\ge1\) and \(\varepsilon>0\), uniformly
over compatible prescribed affine events at \(s\) distinct
sources,
\begin{equation}
 \sum_{z\notin\{x_1,\ldots,x_s\}}
 \P(J_z=1\mid J_{x_1}=\cdots=J_{x_s}=1)
       \ll_{s,C,\varepsilon}N^\varepsilon .
 \label{p1:fixed:conditional}
\end{equation}
If a family of ordered \(s\)-tuples has total positive joint mass \(W_{s,N}\),
the mass of its extensions to ordered distinct \(r\)-tuples,
for fixed \(r\ge s\), is
\begin{equation}
 W_{s\to r,N}\ll_{r,C,\varepsilon}N^\varepsilon W_{s,N}.
 \label{p1:fixed:propagate}
\end{equation}
Both \eqref{p1:fixed:conditional} and \eqref{p1:fixed:propagate}
also hold for homogeneous prescriptions.
For the original run-start events, the mass of fixed \(r\)-tuples
containing a local defect or a compatible contact is
\(O_{r,C,\varepsilon}(N^{-1+\varepsilon})\).
\end{proposition}

\begin{proof}
For \(\delta_z<\eta L\), the conditional sum is at most
\(2Np2^{\eta L}\ll_C N^\eta\). For
\(\delta_z\ge\eta L\), Proposition~\ref{p1:fixed:targets}
counts \(O(N^{\varepsilon/2})\) targets, each of conditional
probability at most one. Take \(\eta=\varepsilon/2\).
The proof uses only the rank loss and compatibility, so
it applies unchanged to homogeneous prescriptions.
For \(r=s\), \eqref{p1:fixed:propagate} is immediate.
For \(r>s\), iterating with margin \(\varepsilon/(r-s)\)
proves the bound. Allowing a seed at any positions
in the ordered tuple multiplies the estimate by a constant
depending on \(r\).

For the original run-start events, the local and contact seeds have positive masses
\(N^{-1+o(1)}\), by the regular-core estimates in
Section~\ref{p1:arith:section}. Contacts at distance
strictly less than \(L\) have zero joint probability.
Propagate the remaining seeds to obtain the last claim.
\end{proof}

\subsection{The correction to every fixed moment}

\begin{theorem}[Fixed factorial moments]
\label{p1:fixed:moments}
For every fixed integer \(r\ge2\) and every
\(\varepsilon>0\),
\begin{equation}
 \E(Z_N)_r
   =\mu_N^r+\binom r2\mu_N^{r-2}\alpha_N
       +O_{r,C,\varepsilon}(N^{-4/5+\varepsilon}).
 \label{p1:fixed:moment-expansion}
\end{equation}
Equivalently, one may use \(\lambda_N\) and
\(a_N=\E(Z_N)_2-\lambda_N^2\).
All fixed ordinary moments are bounded. More precisely,
with \(T_r(u)=\E[\Pois(u)^r]\),
\begin{equation}
 \E Z_N^r
    =T_r(\mu_N)+\frac{\alpha_N}{2}T_r''(\mu_N)
       +O_{r,C,\varepsilon}(N^{-4/5+\varepsilon}).
 \label{p1:fixed:ordinary}
\end{equation}
For every fixed \(r\ge3\), the \(r\)-th factorial
cumulant is \(O_{r,C,\varepsilon}(N^{-4/5+\varepsilon})\).
\end{theorem}

This first estimate is proved independently of the law approximation.
Its bounded-moment conclusion is used in the sharp lower bound of
Theorem~\ref{p13:lower:optimal} in the article.
After that reconstruction, Corollary~\ref{p13:optimal:weighted}
improves the scalar moment remainders to
\(O_{r,C,\varepsilon}(N^{-4/3+\varepsilon})\) when the exact
parameters are retained.

\noindent The complete proof is given in Section~\ref{p14:comp:fixed-moments}.


For \(r\ge2\), the first discrepancy in
\eqref{p1:fixed:moment-expansion} has the exact scale
\(L N^{-2/3}\), by \eqref{p1:eq:alpha-leading}.
The theorem fixes \(r\) first. Its proof does not
allow \(r\) to grow like \(\log N/\log\log N\).
The uniform certificate tools in Section~\ref{p14:cert:section} and the
large-order reconstruction in Section~\ref{p13:transfer:section}
supply that separate input.

\subsection{Fixed-source geometry and the independent moment proof}
\label{p14:comp:fixed}

The statements of the fixed-source estimates and the fixed-moment theorem
are in Section~\ref{p1:fixed:section}. This section supplies the geometric
count and the positive classification used in their proofs. The original
window estimates below use the regular-set convention $p=2^{-L}$;
Section~\ref{p13:exact:section} gives their extension to effective spaces.
The moment proof is independent of the reconstruction of the complete law.

\subsubsection{Uniformly few targets of large rank loss}

We first record the one-parameter geometric estimate in the
precise form needed here.

\begin{lemma}[Several disjoint square equations]
\label{p1:fixed:curve}
Fix \(K,A,D_0>0\) and \(\varepsilon>0\).
There is a fixed integer \(j\) such that the following holds.
Let \(P_1,\ldots,P_j\) be nonconstant products of distinct
linear factors \(Z+u\), of degrees at most \(K\), with
\(|u|\le D_0B\). Suppose their root lists are pairwise disjoint,
and let \(0<|D_\nu|\le N^A\) be integers. Then
\[
 \#\{z\asymp N:
       w_\nu^2=D_\nu P_\nu(z)\ (1\le\nu\le j)
       \text{ for some integers }w_\nu\}
       \ll_{K,A,D_0,\varepsilon}N^\varepsilon .
\]
The constants are uniform in the displayed coefficients and
offsets.
\end{lemma}

\begin{proof}
The square classes of the functions \(D_\nu P_\nu(Z)\) are
geometrically independent: a root of \(P_\nu\) has odd
valuation in that class and even valuation in every other
class. Thus the resulting curve is geometrically integral
and has degree \(2^j\) over the \(Z\)-line.

Choose an integer \(D\ge(K+A)/2\). Give \(T,S\) degree one
and each \(W_\nu\) degree \(D\), and use the equations
\[
 W_\nu^2
   =D_\nu S^{2D-\deg P_\nu}
       \prod_{u\in P_\nu}(T+uS).
\]
The graded ring is free of rank \(2^j\) over \(\mathbb Q[T,S]\);
its generic localization is a field. Its \(D\)-th Veronese
ring is generated in degree one by the monomials in \(T,S\)
of degree \(D\) and the \(W_\nu\). The corresponding
projective curve has degree \(D2^j\), since its hyperplane
bundle is the pullback of \(\mathcal O_{\mathbb P^1}(D)\).
On \(S=1\), the coordinates are integral and have height
\(O_{K,D}(N^D)\).

Pila's uniform curve estimate
\cite[Theorem~A]{pila1995} bounds their number by
\[
 O_{K,A,j,\varepsilon'}
 \bigl((N^D)^{1/(D2^j)+\varepsilon'}\bigr)
       =O(N^{2^{-j}+D\varepsilon'}).
\]
Choose \(j\) and then \(\varepsilon'\) to put this exponent
below \(\varepsilon\). Both choices precede \(N\).
\end{proof}

\begin{lemma}[A product certificate with fixed sources]
\label{p1:fixed:CRT}
Fix \(s,k\ge1\). Uniformly over sources
\(x_1,\ldots,x_s\in I_N\), the number of centres \(z\in I_N\)
for which there are pairwise coprime \(G_1,\ldots,G_k>1\)
satisfying
\[
 G_i\mid x_{a_i}+u_i,\qquad G_i\mid z+w_i,\qquad
 \prod_iG_i>N,
 \quad u_i,w_i\in\{-1,\ldots,L-1\},
\]
is \(O_{s,k,C,\varepsilon}(N^\varepsilon)\) for every
\(\varepsilon>0\).
\end{lemma}

\begin{proof}
There are \((sB)^kB^k=N^{o(1)}\) choices of positions.
For fixed positions, the product \(Q=\prod_iG_i\) divides
\(\prod_i(x_{a_i}+u_i)\le(3N)^k\). Its divisors and their
ordered factorizations into \(k\) factors have
\(N^{o(1)}\) possibilities. Coprimality fixes one residue
class for \(z\) modulo \(Q>N\), hence at most one centre
in \(I_N\). This counts the certificate without requiring
the groups to be fixed in advance.
\end{proof}

\subsubsection{Proof of Proposition~\ref{p1:fixed:targets}}
\label{p14:comp:fixed-targets}

\begin{proof}
Take \(H=D_0B\), with fixed \(D_0>1\). A prime exceeding
\(H\) has at most one occurrence of odd valuation in each
window. Apply Lemma~\ref{p1:fixed:recovery} to
\(R(\boldsymbol x,z)\), with \(d=s\), choosing \(t\) so that
the loss of dimension is at most \(\eta L/4\).
The surviving target projection has dimension at least
\(3\eta L/4\). Nonzero labels have multiplicity at most
\(K=t-1\).

For each prime \(q>H\) incident to an active target coordinate
whose label has odd multiplicity in the \(q\)-row, choose a
minimal circuit of distinct labels containing this label.
Such a circuit exists: a minimal zero-sum subset containing
the label cannot contain a proper zero-sum subset, because
it could be removed. Choose incident coordinates for its
other labels. They all belong to sources, because \(q\)
has only one target occurrence. Group primes according to
the selected physical support. These groups are pairwise
coprime. The number of supports through a fixed target
coordinate is at most
\[
 D=\max\left(1,\sum_{j=2}^{s}j!t^{j-1}K^j\right).
\]

Choose the integer \(j\) in Lemma~\ref{p1:fixed:curve}
for degree at most \(K\) and coefficient height
\(N^{A_0+2K+2}\), where
\(\prod_{q\le H}q\le N^{A_0}\).
Choose \(k\) large enough that
\(\theta=KD/k\le1\) and \(j\theta<\varepsilon/4\).
If the product of the \(k\) largest nontrivial groups
exceeds \(N\), Lemma~\ref{p1:fixed:CRT} applies. It remains
to count the complementary case.

Discard from selection the label classes meeting the target
positions of these \(k\) groups. At most \(k\) classes are
discarded; no further restriction of the relation space is
made. Every group incident to a remaining target class
has size at most \(N^{1/k}\).
For such a class \(A\), write
\[
 M_A=\prod_{v\in A\ {\rm at\ sources}} n_v,\qquad
 P_A(Z)=\prod_{v\in A\ {\rm at\ target}}(Z+w_v).
\]
Its squarefree product \(d_A=\sf(M_AP_A(z))\) can be
factored as
\[
 d_A=d_{0,A}d_{1,A}e_A,\qquad
 d_{0,A}\mid\prod_{q\le H}q,\quad
 d_{1,A}\mid M_A,\quad 1\le e_A\le N^\theta .
\]
Indeed, every remaining prime with odd multiplicity and
not dividing \(M_A\) lies in one of at most \(KD\) incident
groups. The choices of \(d_{0,A}\) and \(d_{1,A}\) are
\(N^{o(1)}\), uniformly for fixed source positions.

There are at least \(3\eta L/4-k\) available classes
meeting the target. Select \(j\) of them. Their target
polynomials have simple disjoint roots. After choosing
their positions and coefficients, the equations
\[
 w_\nu^2=d_{A_\nu}M_{A_\nu}P_{A_\nu}(z)
\]
are counted by Lemma~\ref{p1:fixed:curve}.
The positions cost \(B^{O(1)}\), and the coefficients cost
\(N^{j\theta+o(1)}\). The stated choices, with the error
margin split between these factors, give
\eqref{p1:fixed:targets-bound}.
Only the selected positions and coefficients have been
enumerated, not the possible relation spaces or their
regularizations.
\end{proof}
\subsubsection{Primitive families at a fixed order}

A family is \emph{primitive} if it has a nonzero relative
relation and every proper subfamily has none. Its relation
space projects injectively to each block. For original
windows its nullity is consequently at most \(L\).
For the homogeneous prescription define
\[
 W(I)=p^{|I|}2^{\rho(I)}.
\]
This is an actual joint probability, regardless of the
compatibility of the run-start prescriptions. Those have
joint probability either \(W(I)\) or zero.

\begin{corollary}[Primitive families at a fixed order]
\label{p1:fixed:primitive}
For every fixed $r\ge3$ and $\varepsilon>0$, the sum of the homogeneous
weights $p^r2^{\rho(I)}$ over ordered primitive $r$-tuples of separated
regular windows is $O_{r,C,\varepsilon}(N^{-r/2+\varepsilon})$.
The same bound holds for their actual joint mass and absolute centered mass.
\end{corollary}
\begin{proof}
Apply Proposition~\ref{p13:transfer:primitive} on $G_N$, where
$q_0(I)=p^r2^{\rho_{\mathrm{rel}}(I)}$, and multiply by the fixed
ordering factor $r!$. The proposition uses the geometric half-count and
dense-host construction, independently of the moment proof below.
\end{proof}

The proof is homogeneous throughout. It therefore bounds
incompatible configurations as well as compatible ones.
This feature will be needed for centered Fourier coefficients.

\begin{proposition}[Fixed primitive mass after pair compression]
\label{p1:fixed:compressed}
Fix a channel \(a:b\in\{2:3,3:4\}\), and put
\(\kappa_b=2-1/b\). All original starts lie in \(G_N\) and are
distinct and pairwise separated. For every fixed \(r\ge3\) and
\(\varepsilon>0\), the total
homogeneous mass of primitive families consisting of
one exactly compressed pair in this channel and
\(r-1\) further original windows is
\begin{equation}
 O_{r,C,\varepsilon}
       (N^{-\kappa_b+\varepsilon}).
 \label{p1:fixed:compressed-bound}
\end{equation}
For a compressed geometric good pair and just one further original
window with positive relative rank loss, the mass is
\(O_{C,\varepsilon}(N^{-41/30+\varepsilon})\) in either channel.
The affine positive masses satisfy the same bounds.
\end{proposition}

\begin{proof}
Put $t=r-1\ge2$ and fix a compression motif. The compressed block
is the even-selection image of a genuine set of $O(B)$ values of
width $O(B)$; write $r_e$ for its rank.
Use its normalized centre $bx$; it lies above the original target
range and determines the source pair once the motif is fixed.
There are $B^{O(1)}$ motifs and $NB^{O(1)}$ source pairs.
For its homogeneous mass, Lemma~\ref{p1:arith:alignedrank} gives
$q_0(e)=2^{-r_e}\ll N^{-\kappa_b+\varepsilon'}$, including
affine-incompatible sources. Primitivity gives $\rho\le L$ and
the exact homogeneous weight $q_0(e)p^t2^\rho$.

Every primitive family has a nonzero effective projection in each
of its $t+1$ blocks. Theorem~\ref{p13:half:conditional}, with no old
sources, counts at most $N^{(t+1)/2+\varepsilon'}$ hosts,
including the compressed centre. For $t\ge3$, their total mass is
\[
 N^{(t+1)/2+\varepsilon'}N^{-\kappa_b-t+1+\varepsilon'}
   =N^{-\kappa_b+(3-t)/2+2\varepsilon'}
   \le N^{-\kappa_b+2\varepsilon'}.
\]
For $t=2$, the same count on $\rho\le L/4$ gives
$N^{-\kappa_b-1/4+2\varepsilon'}$.
On $\rho\ge L/4$, fix the compressed source image $W$.
Every proper subfamily is direct, so each target is disjoint from
$W$ and the two targets are direct. Lemma~\ref{p13:transfer:dense-targets}
therefore gives conditional homogeneous target mass
$O(N^{-1+\varepsilon'})$; the original targets are regular, so their
homogeneous marginals are exactly $p$.
The source sum is $\sum_e q_0(e)\ll N^{1-\kappa_b+2\varepsilon'}$,
giving total $O(N^{-\kappa_b+3\varepsilon'})$.
Reduce the margins and absorb the motif choices to obtain
\eqref{p1:fixed:compressed-bound}.

The one-target clause is part (2) of
Proposition~\ref{p1:arith:fixed-homogeneous-seeds}.
Original affine joint probabilities are bounded by the homogeneous ones.
\end{proof}

\subsubsection{Proof of Theorem~\ref{p1:fixed:moments}}
\label{p14:comp:fixed-moments}

\begin{proof}
Expand \(\E(Z_N)_r\) over ordered distinct centres.
We describe all configurations removed before evaluating
the two principal classes.

Local defects and contacts have mass \(N^{-1+\varepsilon}\)
after propagation. A dependent separated pair outside
the two retained channels has mass
\(N^{-4/5+\varepsilon}\), also after propagation.
Tuples containing a primitive original subfamily of
size at least three have that same bound, by
Corollary~\ref{p1:fixed:primitive}.
Two distinct good pairs, including pairs sharing a vertex, have
positive seed mass $O_C(L^2N^{-4/3})$ by
Proposition~\ref{p1:arith:fixed-homogeneous-seeds}(3): their union
belongs to the separated raw-edge population counted there, and
its affine mass is bounded by its homogeneous mass. Fixed-order
propagation preserves the required precision.

There remains the case of one good pair with an
additional relation. Compress this pair exactly.
A minimal dependent subfamily of the compressed
system has at least two blocks and contains the
compressed block, unless an original primitive family
has already been removed. Apply
Proposition~\ref{p1:fixed:compressed} to this seed and
then propagate. Each of its exponents exceeds \(4/5\).
This exhausts the dependent residual configurations:
after their removal the spaces are either all direct,
or one good pair is direct from all remaining
singletons.

The support estimate
of Proposition~\ref{p1:arith:fixed-support}, together with
the \(O_r(N^{r-1}B)\) contact tuples, shows that the
number of full-rank tuples is
\[
 N^r+O_{r,C,\varepsilon}(N^{r-4/5+\varepsilon}).
\]
Their joint probability is exactly \(p^r\), and their
contribution is \(\lambda_N^r+O(N^{-4/5+\varepsilon})\).

Let \(A_N\) be the ordered positive mass of good pairs.
The pair calculation gives
\[
 A_N=a_N+O_{C,\varepsilon}(N^{-4/5+\varepsilon}).
\]
Fix one such pair. After any fixed number of already
chosen windows, at most \(N^{1/2+o(1)}\) target centres
are dependent on their relative space, by
Proposition~\ref{p1:arith:fixed-branching}. Thus the number of
ordered direct extensions by \(r-2\) windows is
\[
 N^{r-2}\bigl(1+O_{r,C}(N^{-1/2+o(1)})\bigr),
\]
uniformly in the pair. Their joint weight is
\(q(e)p^{r-2}\). Placing the ordered pair in two of
the \(r\) positions gives contribution
\[
 \binom r2\lambda_N^{r-2}A_N
       +O_{r,C}(A_NN^{-1/2+o(1)}).
\]
Since \(A_N=O_C(LN^{-2/3})\), the last error is
smaller than the stated remainder. The two principal
classes give the formula with \((\lambda_N,a_N)\).
The mean transfer
\(\mu_N-\lambda_N=O(N^{-1+\varepsilon})\) and
\(\alpha_N-a_N=O(N^{-1+\varepsilon})\) give
\eqref{p1:fixed:moment-expansion}.

The Stirling-number identity
\(Z^r=\sum_{j=0}^r\left\{\begin{smallmatrix}r\\j\end{smallmatrix}\right\}(Z)_j\)
proves \eqref{p1:fixed:ordinary} and boundedness.
Finally, to any fixed formal order, the factorial
generating function equals
\[
 e^{\mu_N t}\left(1+\frac{\alpha_Nt^2}{2}\right)
       +O_{\rm coeff}(N^{-4/5+\varepsilon}).
\]
Taking its formal logarithm is legitimate at that
fixed order. Terms quadratic in \(\alpha_N\) are
\(O(N^{-4/3+o(1)})\), which proves the cumulant claim.
\end{proof}
